%=====================================================================
\chapter{Brute Force, and When It Is the Right Answer}\label{ch:bruteforce}
%=====================================================================
\locator{2}{Part II --- Recursion and Divide-and-Conquer}

\begin{outcomes}
By the end of this chapter you will be able to:
\begin{tight}
  \item \textbf{Define} the brute-force strategy and \textbf{apply} it to a
        stated problem \emph{(CLO-6)}.
  \item \textbf{Distinguish} removing redundant work within a search from
        changing the search itself, and \textbf{say} which one changes the
        order of growth.
  \item \textbf{Derive} the running time of the three maximum-subarray
        algorithms and \textbf{verify} the predicted ratios by measurement.
  \item \textbf{State} the conditions under which brute force is the correct
        engineering choice.
  \item \textbf{Use} a brute-force implementation as a correctness oracle for a
        cleverer one.
\end{tight}
\end{outcomes}

\begin{prereq}
\chref{asymptotics} for the notation, \chref{invariants} for the invariant in
Section~7.3. No recursion is needed; this chapter is the baseline the rest of
Part~II improves on.
\end{prereq}

\begin{keyterms}
brute force $\bullet$ exhaustive search $\bullet$ candidate solution $\bullet$
maximum subarray $\bullet$ redundant work $\bullet$ incremental computation
$\bullet$ Kadane's algorithm $\bullet$ oracle $\bullet$ baseline
\end{keyterms}

\section{The Strategy With No Ideas In It}

\begin{definitionbox}[Brute Force]
\term{Brute force} --- also \term{exhaustive search} --- solves a problem by
generating every candidate solution and testing each one. It needs no insight
into the problem's structure beyond the ability to recognise an answer.

\smallskip
It is the first algorithm you should write and, often, the last one you should
ship.
\end{definitionbox}

\begin{conceptbox}[Why a chapter on the obvious]
The published HEC outline lists brute force among the design techniques, beside
divide-and-conquer, greedy and dynamic programming. That is the right company
for it, for three reasons.

\begin{tightnum}
  \item \textbf{It is the baseline.} ``Quicksort is fast'' is meaningless
        without something to be fast relative to. Every later chapter is an
        improvement on an exhaustive search, and you cannot measure an
        improvement against nothing.
  \item \textbf{It is the oracle.} A brute-force implementation is short enough
        to be obviously correct, so it is what you test the clever algorithm
        \emph{against}. Section~7.4 does exactly this.
  \item \textbf{It is sometimes the answer.} For small or bounded $n$, for code
        that runs once, or for a correctness-critical path, the four-line
        version you can read in one sitting beats the forty-line version you
        cannot.
\end{tightnum}
\end{conceptbox}

\section{A Problem With Three Answers}

\begin{definitionbox}[Maximum Subarray]
Given an array $A[1 \mathinner{\ldotp\ldotp} n]$ of numbers, some negative,
find the contiguous block $A[i \mathinner{\ldotp\ldotp} j]$ with the largest
sum.

\smallskip
If all the numbers were positive the answer would be the whole array and there
would be no problem. The negatives are what make it one.
\end{definitionbox}

\dsfig{\aoaalg{\algname{Max-Subarray-Brute}$(A, n)$}{
\State $\mathit{best} \gets -\infty$
\For{$i \gets 1$ \textbf{to} $n$}
  \For{$j \gets i$ \textbf{to} $n$}
    \State $s \gets 0$
    \For{$k \gets i$ \textbf{to} $j$}
      \State $s \gets s + A[k]$
        \Comment{re-added from scratch}
    \EndFor
    \State $\mathit{best} \gets \max(\mathit{best}, s)$
  \EndFor
\EndFor
\State \Return $\mathit{best}$
}}{\algname{Max-Subarray-Brute}: every block, summed from scratch. Three nested
loops, so $\bigTh{n^{3}}$ --- and lines 4 to 7 are doing work that lines 4 to 7
of the previous iteration already did.}{brute3}

\begin{theorem}[The brute force is cubic]\label{thm:brute3}
\algname{Max-Subarray-Brute} runs in $\bigTh{n^{3}}$ time.
\end{theorem}

\begin{proof}
For each pair $(i,j)$ with $i \le j$, the inner loop performs $j - i + 1$
additions. Summing over all pairs,
\[ \sum_{i=1}^{n}\sum_{j=i}^{n} (j-i+1)
   = \sum_{i=1}^{n}\sum_{m=1}^{n-i+1} m
   = \sum_{i=1}^{n} \frac{(n-i+1)(n-i+2)}{2}
   = \sum_{p=1}^{n} \frac{p(p+1)}{2}. \]
The last sum is $\tfrac12\bigl(\sum p^{2} + \sum p\bigr) =
\tfrac12\bigl(\tfrac{n(n+1)(2n+1)}{6} + \tfrac{n(n+1)}{2}\bigr)$, whose leading
term is $n^{3}/6$. By Lemma~\ref{lem:poly} this is $\bigTh{n^{3}}$. Every other
operation is $\bigO{1}$ per pair, contributing $\bigO{n^{2}}$, which the sum
rule absorbs.
\end{proof}

\subsection*{Removing the waste, which is not a new idea}

The inner loop recomputes a sum the previous iteration almost had. Extending
the block from $j$ to $j+1$ costs one addition.

\dsfig{\aoaalg{\algname{Max-Subarray-Quadratic}$(A, n)$}{
\State $\mathit{best} \gets -\infty$
\For{$i \gets 1$ \textbf{to} $n$}
  \State $s \gets 0$
  \For{$j \gets i$ \textbf{to} $n$}
    \State $s \gets s + A[j]$
      \Comment{carried, not rebuilt}
    \State $\mathit{best} \gets \max(\mathit{best}, s)$
  \EndFor
\EndFor
\State \Return $\mathit{best}$
}}{\algname{Max-Subarray-Quadratic}. The search is unchanged --- still every
pair $(i,j)$ --- but each candidate now costs $\bigO{1}$ instead of
$\bigO{n}$. A better accountant, not a better idea.}{brute2}

\begin{alertbox}[This is the distinction the chapter exists for]
Going from $n^{3}$ to $n^{2}$ removed \emph{redundant work inside} the search.
The set of candidates examined did not change: still all $\binom{n}{2}$ pairs.

\smallskip
Going from $n^{2}$ to $n$, next, will change \emph{the search itself} --- it
will stop examining pairs at all.

\smallskip
Both are worth doing and they are different activities. The first is careful
programming and gets you a constant or a factor of $n$; the second requires an
idea about the problem, and it is what the rest of this book is about.
\end{alertbox}

\section{An Idea About the Problem}

\begin{conceptbox}[Kadane's observation]
Ask a smaller question: \emph{what is the best block \textbf{ending exactly at
position $j$}?}

\smallskip
Such a block either extends the best block ending at $j-1$, or it starts fresh
at $j$. There is no third option --- a block ending at $j$ that does not
contain $j-1$ must begin at $j$. So if $H_{j}$ is the best sum ending at $j$,
\[ H_{j} = \max\bigl(A[j],\; H_{j-1} + A[j]\bigr), \]
and the answer is $\max_{j} H_{j}$.

\smallskip
That one sentence removes an entire dimension from the search. It is also the
first appearance in this book of \emph{optimal substructure} --- the best
solution built from best solutions to smaller versions --- which \chref{dpone}
makes into a general technique.
\end{conceptbox}

\dsfig{\aoaalg{\algname{Max-Subarray-Kadane}$(A, n)$}{
\State $\mathit{here} \gets 0$;\quad $\mathit{best} \gets -\infty$
\For{$j \gets 1$ \textbf{to} $n$}
  \State $\mathit{here} \gets \max\bigl(A[j],\, \mathit{here} + A[j]\bigr)$
    \Comment{extend, or start fresh}
  \State $\mathit{best} \gets \max(\mathit{best}, \mathit{here})$
\EndFor
\State \Return $\mathit{best}$
}}{\algname{Max-Subarray-Kadane}. One pass, two variables. The loop invariant
is what makes it believable, and Theorem~\ref{thm:kadane} is what makes it
true.}{kadane}

\begin{theorem}[Kadane's algorithm is correct]\label{thm:kadane}
\algname{Max-Subarray-Kadane} returns the maximum subarray sum, in
$\bigTh{n}$ time.
\end{theorem}

\begin{proof}
Invariant: \emph{at the start of the iteration for $j$, $\mathit{here}$ is the
maximum sum of a block ending at $j-1$, and $\mathit{best}$ is the maximum sum
of any block within $A[1 \mathinner{\ldotp\ldotp} j-1]$.}

\smallskip
\textbf{Initialisation.} Before $j = 1$ both ranges are empty;
$\mathit{best} = -\infty$ correctly records that no block has been seen, and
$\mathit{here} = 0$ is the empty sum.

\smallskip
\textbf{Maintenance.} Assume the invariant before iteration $j$. Any block
ending at $j$ either contains $j-1$ --- in which case its sum is (the sum of a
block ending at $j-1$) $+\, A[j]$, maximised by $\mathit{here} + A[j]$ using
the invariant --- or it does not, in which case it is $A[j]$ alone. Line 3
takes the larger of exactly these two, so afterwards $\mathit{here}$ is the
best sum ending at $j$. Line 4 then extends the running maximum over one more
position, re-establishing the invariant for $j+1$.

\smallskip
\textbf{Termination.} The loop runs $n$ times and stops. Substituting $j = n+1$
into the invariant, $\mathit{best}$ is the maximum sum of any block within
$A[1\mathinner{\ldotp\ldotp}n]$, which is the postcondition. The body is
$\bigO{1}$ and runs $n$ times, so the time is $\bigTh{n}$.
\end{proof}

\section{The Three, Raced}

\begin{worked}[First, Do They Agree?]
A faster algorithm that returns a different answer is not faster. Before any
timing, \texttt{code/ch07\_brute.cpp} runs all three on the same inputs:

\rowcolors{2}{white}{lightbg}
\begin{center}
\begin{tabular}{@{}rrrrl@{}}
\toprule
\rowcolor{primary!15}
$n$ & cubic & quadratic & linear & \\
\midrule
  1 &   35 &   35 &   35 & agree \\
  2 &   80 &   80 &   80 & agree \\
 10 &  265 &  265 &  265 & agree \\
100 &  563 &  563 &  563 & agree \\
500 & 1362 & 1362 & 1362 & agree \\
\bottomrule
\end{tabular}
\end{center}

\smallskip
This is the oracle use of brute force, and it is not a formality. The cubic
version is three nested loops with no ideas in it and is almost impossible to
get wrong; Kadane's is four lines that are easy to get subtly wrong --- initialise
$\mathit{best}$ to $0$ instead of $-\infty$ and it silently returns $0$ for an
all-negative array. The agreement above is evidence, cheaply bought.
\end{worked}

\begin{worked}[The Predicted Ratios: 8, 4, 2]
Doubling $n$ should multiply a cubic time by 8, a quadratic by 4 and a linear
by 2. Median of 3, 15 and 101 runs respectively; times in milliseconds.

\rowcolors{2}{white}{lightbg}
\begin{center}
\begin{tabular}{@{}rrrrrrr@{}}
\toprule
\rowcolor{primary!15}
$n$ & cubic & ratio & quadratic & ratio & linear & ratio \\
\midrule
   500 &    7.66 & ---  & 0.070 & ---  & 0.0005 & ---  \\
1{,}000 &   60.25 & 7.86 & 0.304 & 4.32 & 0.0010 & 2.00 \\
2{,}000 &  471.16 & 7.82 & 1.095 & 3.60 & 0.0018 & 1.80 \\
4{,}000 & 3355.85 & 7.12 & 4.348 & 3.97 & 0.0035 & 1.94 \\
8{,}000 & ---     & ---  & 17.299 & 3.98 & 0.0071 & 2.03 \\
\bottomrule
\end{tabular}
\end{center}

\smallskip
\textbf{Means: 7.60, 3.97, 1.94} against predicted 8, 4, 2.

\smallskip
\textbf{Now read the first and last columns together.} At $n = 4{,}000$ the
cubic version takes 3.36 \emph{seconds} and the linear version takes 3.5
\emph{microseconds} --- a factor of about a million. At $n = 500$ the factor was
about fifteen thousand. The gap is not a constant; it grows as $n^{2}$, which
is exactly what ``$\bigTh{n^{3}}$ against $\bigTh{n}$'' means in seconds.

\smallskip
\textbf{And the row that is missing} --- cubic at $n = 8{,}000$ --- is missing
because it would take about 27 seconds for a result the linear version produces
in 7 microseconds.
\end{worked}

\section{When Brute Force Is the Right Answer}

The table above makes brute force look indefensible. In absolute terms it often
is not.

\begin{worked}[What It Costs, in Absolute Terms]
The same cubic routine, read as a cost rather than as a ratio:

\rowcolors{2}{white}{lightbg}
\begin{center}
\begin{tabular}{@{}rrl@{}}
\toprule
\rowcolor{primary!15}
$n$ & cubic & verdict \\
\midrule
   10 &    0.1 $\mu$s & imperceptible \\
   30 &    2.1 $\mu$s & imperceptible \\
  100 &   78.8 $\mu$s & imperceptible \\
  300 &    1.6 ms     & a visible pause \\
1{,}000 &   52.3 ms   & a visible pause \\
3{,}000 & 1470.1 ms   & too slow to ship \\
\bottomrule
\end{tabular}
\end{center}

\smallskip
\textbf{At $n = 100$ the cubic algorithm takes 79 microseconds.} If your arrays
are bounded by a hundred elements and the routine runs once per user action,
the asymptotically terrible algorithm is the correct engineering choice: it is
four lines, it needs no proof, and nobody will ever notice.

\smallskip
\textbf{The decision turns on the range of $n$, not on the exponent.} Which is
why the first question to ask about any algorithm is not ``what is its order of
growth'' but ``how big does $n$ get''.
\end{worked}

\begin{conceptbox}[A checklist]
Brute force is defensible when \emph{all} of these hold:

\begin{tight}
  \item $n$ is small and \textbf{bounded by something real} --- a board size, a
        number of columns, a field count --- not merely small today;
  \item the code runs rarely, or off the critical path;
  \item correctness matters more than speed, and the short version is the one
        you can convince yourself of;
  \item you have measured, at the largest $n$ you will actually see.
\end{tight}

\smallskip
It is indefensible when $n$ is unbounded or user-supplied, because then you
have not chosen a slow algorithm, you have chosen an unbounded one --- and
\chref{complexity} is about problems where that is the only choice anyone
knows.
\end{conceptbox}

\begin{pitfall}
\begin{tight}
  \item \textbf{Confusing ``removed redundant work'' with ``found a better
        algorithm''.} The cubic-to-quadratic step did not change which
        candidates were examined. The quadratic-to-linear step did.
  \item \textbf{Initialising $\mathit{best}$ to 0 in Kadane's algorithm.} It
        then returns 0 on an all-negative array, which is wrong unless the
        empty block is allowed --- and if it is, say so in the specification.
  \item \textbf{Dismissing brute force without measuring.} 79 microseconds at
        $n = 100$ is not a performance problem.
  \item \textbf{Keeping brute force when $n$ is user-supplied.} ``Small in
        practice'' is a prediction, and users falsify it.
  \item \textbf{Not writing the brute force at all.} You then have nothing to
        test the clever version against, which is how subtly wrong algorithms
        reach production.
  \item \textbf{Timing the brute force once and the clever one once.} At
        microsecond scale one run measures the scheduler.
\end{tight}
\end{pitfall}

\begin{lab}[Three Algorithms, One Problem]
Reproduces all three tables. Expect twenty-five minutes; the cubic column at
$n = 4{,}000$ takes about ten seconds on its own.

\begin{lstlisting}[language=C++]
// ch07_brute.cpp -- brute force, and the two things that beat it.
// Build: cl /O2 /EHsc /std:c++17 ch07_brute.cpp     (see Appendix A)

// --- 1. Brute force, as literally as it can be written -------------
// For every start i and end j, add up a[i..j] from scratch. The inner
// re-addition is pure waste, and it is the n^3.
static long long max_sub_cubic(const std::vector<int>& a) {
    long long best = std::numeric_limits<long long>::min();
    size_t n = a.size();
    for (size_t i = 0; i < n; ++i)
        for (size_t j = i; j < n; ++j) {
            long long s = 0;
            for (size_t k = i; k <= j; ++k) s += a[k];
            best = std::max(best, s);
        }
    return best;
}

// --- 2. The same search, with the waste removed --------------------
// Extending the block from j to j+1 costs one addition, not a fresh
// traversal. The SEARCH is unchanged -- still every (i, j) pair -- so
// this is not yet a better idea, only a better accountant.
static long long max_sub_quadratic(const std::vector<int>& a) {
    long long best = std::numeric_limits<long long>::min();
    for (size_t i = 0; i < a.size(); ++i) {
        long long s = 0;
        for (size_t j = i; j < a.size(); ++j) {
            s += a[j];
            best = std::max(best, s);
        }
    }
    return best;
}

// --- 3. Kadane: a different idea, not a faster loop ----------------
// Invariant: after processing a[0..j], "here" is the largest sum of a
// block ENDING at j, and "best" the largest sum anywhere in a[0..j].
// A block ending at j+1 either extends the best block ending at j, or
// starts fresh. That one observation removes a whole dimension.
static long long max_sub_linear(const std::vector<int>& a) {
    long long best = std::numeric_limits<long long>::min(), here = 0;
    for (int x : a) {
        here = std::max((long long)x, here + x);
        best = std::max(best, here);
    }
    return best;
}

// --- 4. All three agree, which is what makes the race fair ---------
// A faster algorithm that gets a different answer is not faster.

// --- 5. The ratios each order of growth predicts -------------------
// Doubling n should multiply the times by 8, 4 and 2.

// --- 6. Where brute force is the right answer ----------------------
// Report the ABSOLUTE cost, not a ratio against a linear time too
// small to measure, and say where it crosses thresholds a person
// would actually notice.
\end{lstlisting}

\textbf{What to notice.} Part A costs almost nothing to run and is the most
valuable part of the program. Delete it and the other two parts are a race
between three functions you have not checked.

\smallskip
\textbf{Then try to break it.} Change \texttt{best} in
\texttt{max\_sub\_linear} from \texttt{numeric\_limits::min()} to \texttt{0},
and run part A again. Every row still agrees --- because the random inputs all
contain a positive block. Now generate an all-negative array and watch part A
catch it. An oracle only tests the inputs you give it, which is
\chref{cases}'s point about distributions arriving again in a new costume.
\end{lab}

\begin{checkpoint}
\begin{tightnum}
  \item The cubic algorithm performs $\sum_{i}\sum_{j\ge i}(j-i+1)$ additions.
        Evaluate this sum for $n = 100$ and check it against $n^{3}/6$.
  \item Explain why going from $n^{3}$ to $n^{2}$ did not require any insight
        into the maximum-subarray problem, while going from $n^{2}$ to $n$ did.
  \item Give a concrete situation in which you would ship the $\bigTh{n^{3}}$
        version, and state the one fact that would change your mind.
\end{tightnum}
\end{checkpoint}

\begin{chaptersummary}
\begin{tight}
  \item Brute force generates every candidate and tests it. It needs no insight
        into the problem, which is both its weakness and its value.
  \item It earns its place as the \emph{baseline} to improve on and the
        \emph{oracle} to test against --- and sometimes as the answer.
  \item Maximum subarray has three solutions: $\bigTh{n^{3}}$ exhaustive,
        $\bigTh{n^{2}}$ with the redundant re-summation removed, and
        $\bigTh{n}$ by Kadane's observation.
  \item Removing redundant work \emph{within} a search is different from
        changing \emph{the search}. Only the second needs an idea about the
        problem, and it is the one this book is mostly about.
  \item Measured ratios on doubling: 7.60, 3.97, 1.94 against predicted 8, 4,
        2. At $n = 4{,}000$ the cubic version took 3.4 seconds and the linear
        one 3.5 microseconds.
  \item Measured absolutely, the cubic version costs 79 $\mu$s at $n = 100$ and
        1.5 s at $n = 3{,}000$. For a bounded small $n$ it is the right choice.
  \item The first question about an algorithm is how big $n$ gets, not what its
        exponent is.
\end{tight}
\end{chaptersummary}

\begin{reviewq}
  \item \textbf{(MCQ)} The brute-force maximum-subarray algorithm runs in:
        (a)~$\bigTh{n}$; (b)~$\bigTh{n\lg n}$; (c)~$\bigTh{n^{2}}$;
        (d)~$\bigTh{n^{3}}$.
  \item \textbf{(MCQ)} Kadane's algorithm improves on the quadratic version by:
        (a)~using a faster inner loop; (b)~examining fewer candidate blocks;
        (c)~using recursion; (d)~sorting the input first.
  \item \textbf{(Short)} Define the brute-force strategy and give two reasons a
        course on algorithm design includes it \emph{(CLO-6)}.
  \item \textbf{(Short)} State Kadane's loop invariant and explain why only two
        cases need be considered in its maintenance step.
  \item \textbf{(Short)} Explain the oracle use of a brute-force
        implementation, and its one limitation.
  \item \textbf{(Applied)} Give a brute-force algorithm for finding the closest
        pair among $n$ points in the plane, state its running time with
        justification, and say what you would need to beat it.
  \item \textbf{(Applied)} Your team has a $\bigTh{n^{3}}$ routine running on
        $n \le 80$, taking about 40 $\mu$s. A colleague proposes replacing it
        with a $\bigTh{n}$ version, 60 lines long. Argue both sides, and state
        the single piece of information that should decide it.
\end{reviewq}
